% =============================================================================
% Wprowadzenie. Plik edytowany ręcznie -- skrypt analiza_odpowiedzi_naukowa.py go NIE nadpisuje.
% Liczby: definicje.tex; tabele: \TabDir (CSV); rysunki: \FigDir; literatura: references.bib.
% =============================================================================
\section{Wprowadzenie}

Przypominanie sobie informacji z pamięci wzmacnia ich późniejsze utrwalenie silniej niż ponowne
ich przeczytanie (efekt testowania), a nauka rozłożona w czasie jest zapamiętywana lepiej niż nauka
skumulowana (efekt rozłożenia w czasie) \citep{roediger2006,cepeda2006,cepeda2008}. W randomizowanych
badaniach z udziałem lekarzy w trakcie kształcenia tzw. edukacja rozłożona w czasie (ang.
\emph{spaced education}), w której krótkie pytania wielokrotnego wyboru przesyła się cyklicznie
przez tygodnie lub miesiące, poprawiała przyswajanie i utrwalanie wiedzy
\citep{kerfoot2007urology,kerfoot2007mededuc,kerfoot2009}, natomiast przeglądy systematyczne wskazują na małe
lub umiarkowane korzyści u pracowników ochrony zdrowia \citep{phillips2019,martinengo2024}.

W szkoleniach organizacyjnych stosuje się krótkie formy \emph{mikrouczenia} (ang.
\emph{microlearning}) oparte na quizach, jednak dostępne dowody pochodzą głównie z badań na małych
próbach lub z tzw. szarej literatury \citep{taylor2022}. Ewaluacje obowiązkowych szkoleń organizacyjnych, np.\
z zakresu świadomości bezpieczeństwa, często wykazują efekty niewielkie albo zanikające z upływem czasu
\citep{reinheimer2020,lain2022,ho2025}. Dane z dzienników (logów) takich platform pozwalają
natomiast obserwować uczenie się na poziomie pojedynczych odpowiedzi, podobnie jak w modelach ucznia
rozwijanych dla inteligentnych systemów tutorskich i aplikacji do nauki języków
\citep{corbett1994,pavlik2009,settles2016,tabibian2019,koedinger2023}.

Badanie przeprowadzono w banku, w którym pracownicy odpowiadali na krótkie pytania wielokrotnego
wyboru dotyczące wymogów regulacyjnych, produktów bankowych, umiejętności miękkich oraz, w przypadku
kadry kierowniczej, zarządzania, wyświetlane przez aplikację na komputerze w trakcie dnia pracy.
Trzyletnie logi odpowiedzi z tego programu posłużyły do weryfikacji następujących pytań
badawczych:
(PB1) czy odsetek poprawnych odpowiedzi rośnie wraz z kolejnymi przypomnieniami tego samego pytania
przez tego samego pracownika;
(PB2) jak poprawność zależy od wyniku poprzedniej próby i od czasu, który od niej upłynął;
(PB3) czy powtarzane błędy są trwałe, co wskazywałoby raczej na błędne przekonania niż na
zgadywanie;
(PB4) czy poprawa przenosi się na nowo napotykane pytania; oraz
(PB5) jakie są właściwości psychometryczne pytań?
